\section*{Теория}
\addcontentsline{toc}{section}{Теория}

\subsection*{0. Базовая теория}

\noindent\textbf{Пример нахождения ФСР однородной системы $Ax = 0$}

\noindent\textbf{Условие задачи:}
Найти фундаментальную систему решений (ФСР) и общее решение однородной системы линейных уравнений, заданной матрицей коэффициентов $A$:
\[
A = \begin{pmatrix}
2 & -1 & 3 & -7 \\
2 & -1 & 11 & -24 \\
4 & -2 & 14 & -31
\end{pmatrix}
\]

\noindent\textbf{Шаг 1. Приведение матрицы к ступенчатому виду (Метод Гаусса)}
Выполним элементарные преобразования над строками матрицы $A$:
1. Из второй строки вычтем первую ($II - I$).
2. Из третьей строки вычтем две первых ($III - 2\cdot I$).

\[
\begin{pmatrix}
2 & -1 & 3 & -7 \\
2 & -1 & 11 & -24 \\
4 & -2 & 14 & -31
\end{pmatrix}
\xrightarrow{\substack{II - I \\ III - 2\cdot I}}
\begin{pmatrix}
2 & -1 & 3 & -7 \\
0 &  0 & 8 & -17 \\
0 &  0 & 8 & -17
\end{pmatrix}
\xrightarrow{III - II}
\begin{pmatrix}
2 & -1 & 3 & -7 \\
0 &  0 & 8 & -17 \\
0 &  0 & 0 &  0
\end{pmatrix}
\]

Матрица приведена к ступенчатому виду. Количество ненулевых строк определяет ранг матрицы: $\mathrm{Rg}\,A = 2$. \\
Количество переменных (размерность пространства) $n = 4$. \\
Количество векторов в ФСР (число свободных переменных): $K = n - \mathrm{Rg}\,A = 4 - 2 = 2$.

\noindent\textbf{Шаг 2. Разделение переменных на базисные и свободные}
Ступеньки матрицы (ведущие элементы) находятся в столбцах $1$ и $3$. 
\textbf{Базисные (главные) переменные:} $x_1, x_3$.
\textbf{Свободные переменные:} $x_2, x_4$.

\noindent\textbf{Шаг 3. Нахождение векторов ФСР ($f_1$ и $f_2$)}
Чтобы найти линейно независимые векторы ФСР, мы поочередно присваиваем одной из свободных переменных значение единицы, обнуляя остальные свободные переменные.

\noindent
\textbf{Заполнение таблицы значений свободных переменных:}
\[
\begin{array}{|c|c|c|c|c|}
\hline
    & x_1 & x_2 & x_3 & x_4 \\ \hline
f_1 &     &  1  &     &  0  \\ \hline
f_2 &     &  0  &     &  1  \\ \hline
\end{array}
\]

\noindent
\textbf{Вычисление базисных переменных для каждого вектора:}
\begin{itemize}
    \item \textbf{Для первого вектора $f_1$ (подставляем $x_2 = 1, x_4 = 0$):}
    \[ 
    8x_3 - 17(0) = 0 \implies x_3 = 0 
    \]
    \[ 
    2x_1 - 1(1) + 3(0) - 7(0) = 0 \implies 2x_1 = 1 \implies x_1 = \frac{1}{2} 
    \]
    Получаем вектор: $f_1' = \begin{pmatrix} \frac{1}{2} & 1 & 0 & 0 \end{pmatrix}^T$. Чтобы избавиться от дробей, умножим его на $2$:
    \[
    f_1 = \begin{pmatrix} 1 \\ 2 \\ 0 \\ 0 \end{pmatrix}
    \]
    
    \item \textbf{Для второго вектора $f_2$ (подставляем $x_2 = 0, x_4 = 1$):}
    \[ 
    8x_3 - 17(1) = 0 \implies x_3 = \frac{17}{8} 
    \]
    \[ 
    2x_1 - 1(0) + 3\left(\frac{17}{8}\right) - 7(1) = 0 \implies 2x_1 + \frac{51}{8} - \frac{56}{8} = 0 \implies 2x_1 = \frac{5}{8} \implies x_1 = \frac{5}{16} 
    \]
    Получаем вектор: $f_2' = \begin{pmatrix} \frac{5}{16} & 0 & \frac{17}{8} & 1 \end{pmatrix}^T$. Чтобы избавиться от дробей, умножим его на $16$:
    \[
    f_2 = \begin{pmatrix} 5 \\ 0 \\ 34 \\ 16 \end{pmatrix}
    \]
\end{itemize}

\noindent
\textbf{Итоговая таблица векторов ФСР (базис пространства решений):}
\[
\begin{array}{|c|c|c|c|c|}
\hline
    & x_1 & x_2 & x_3 & x_4 \\ \hline
f_1 &  1  &  2  &  0  &  0  \\ \hline
f_2 &  5  &  0  &  34 &  16 \\ \hline
\end{array}
\]

\noindent
\textbf{Ответ: Общее решение ОСЛАУ}
Любое решение системы выражается как линейная комбинация векторов ФСР:
\[
x = \alpha \begin{pmatrix} 1 \\ 2 \\ 0 \\ 0 \end{pmatrix} 
+ \beta \begin{pmatrix} 5 \\ 0 \\ 34 \\ 16 \end{pmatrix}, 
\quad \alpha, \beta \in \mathbb{R}
\]

\begin{definition}
\textit{Кольцом} называется множество $R$, на котором определены бинарные операции «сложения» $+$ и «умножения» $\cdot$, которые удовлетворяют следующим свойствам:

\begin{enumerate}[label=\arabic*)]
    \setcounter{enumi}{-1} % Начинаем нумерацию с нуля
    \item $R$ замкнуто относительно операций сложения и умножения.
    \item $(R, +)$ — абелева группа (коммутативная группа).
    \item $(R, \cdot)$ — полугруппа (ассоциативность умножения).
    \item $\forall a, b, c \in R : a \cdot (b + c) = a \cdot b + a \cdot c$ и $(a + b) \cdot c = a \cdot c + b \cdot c$ (дистрибутивность).
\end{enumerate}
\end{definition}

\begin{definition}
    \textit{Полем} называется коммутативное кольцо $R$ с единицей не равной нулю, в котором каждый ненулевой элемент обратим.
\end{definition}

\begin{definition}
    Пусть $\mathbb{F}$ — поле, $e_1, \dots, e_n \in \mathbb{F}^k$. Будем говорить, что вектора $e_1, \dots, e_n$ \textit{линейно зависимы} над $\mathbb{F}$, если $\exists \lambda_1, \dots, \lambda_n$ такие, что нетривиальная (т.е. $(\lambda_1, \dots, \lambda_n) \neq (0, \dots, 0)$) линейная комбинация $\lambda_1 e_1 + \dots + \lambda_n e_n$ равна нулю. В противном случае будем говорить, что вектора \textit{линейно независимы} над $\mathbb{F}$.
\end{definition}

\begin{definition}
    \textbf{Векторное (линейное) пространство} над полем $\mathbb{F}$ (обычно $\mathbb{F} = \mathbb{R}$ или $\mathbb{C}$) — это множество $V$, элементы которого называются \textbf{векторами}, снабжённое двумя операциями: сложением векторов и умножением вектора на скаляр (элемент поля). Эти операции должны удовлетворять следующим аксиомам.

\textit{Аксиомы сложения:}

\begin{enumerate}
    \item \textbf{Коммутативность:} $x + y = y + x$ для любых $x, y \in V$.
    \item \textbf{Ассоциативность:} $(x + y) + z = x + (y + z)$ для любых $x, y, z \in V$.
    \item \textbf{Существование нулевого вектора:} существует элемент $0 \in V$ такой, что $x + 0 = x$ для любого $x \in V$.
    \item \textbf{Существование противоположного вектора:} для каждого $x \in V$ существует элемент $-x \in V$ такой, что $x + (-x) = 0$.
\end{enumerate}

\textit{Аксиомы умножения на скаляр:}

\begin{enumerate}
    \item \textbf{Дистрибутивность умножения на скаляр относительно сложения векторов:} $\lambda(x + y) = \lambda x + \lambda y$ для любых $\lambda \in \mathbb{F}, x, y \in V$.
    \item \textbf{Дистрибутивность относительно сложения скаляров:} $(\lambda + \mu)x = \lambda x + \mu x$ для любых $\lambda, \mu \in \mathbb{F}, x \in V$.
    \item \textbf{Ассоциативность умножения на скаляр:} $(\lambda\mu)x = \lambda(\mu x)$ для любых $\lambda, \mu \in \mathbb{F}, x \in V$.
    \item \textbf{Унитарность:} $1 \cdot x = x$ для любого $x \in V$, где $1$ — единица поля $\mathbb{F}$.
\end{enumerate}
\end{definition}

\begin{definition}
    Множество векторов $e_1, \dots, e_n$ линейного пространства $L$ является его \textit{базисом}, если:
\begin{enumerate}
    \item[1)] Любой вектор $x \in L$ выражается в виде линейной комбинации векторов $e_1, \dots, e_n$;
    \item[2)] Вектора $e_1, \dots, e_n$ линейно независимы.
\end{enumerate}
\end{definition}

Базис векторного пространства $V$ может быть выбран разными способами, но количество векторов в любом базисе $V$ одинаково. Это количество называется \textbf{размерностью} векторного пространства $V$ и обозначается $\dim V$. Если это число конечно, пространство называется \textbf{конечномерным}.

\begin{definition}
\textbf{Линейный оператор} (или линейное преобразование) пространства $V$ — это отображение $\varphi : V \to V$, удовлетворяющее для любых векторов $x, y \in V$ и любого скаляра $\lambda \in \mathbb{F}$ условиям:
$$ \varphi(x + y) = \varphi(x) + \varphi(y), \quad \varphi(\lambda x) = \lambda\varphi(x). $$
\end{definition}

\begin{definition}
    Подмножество $M$ линейного пространства $L$ над полем $\mathbb{F}$ называется \textit{линейным подпространством}, если:
\begin{enumerate}
    \item[1)] Для любых $x \in M$ и $\alpha \in \mathbb{F}$: $\alpha x \in M$;
    \item[2)] Для любых $x, y \in M$: $x + y \in M$.
\end{enumerate}

\end{definition}

\subsection*{1. Алгебра над полем}

\begin{definition}[Э.Б.Винберг. Курс алгебры -- стр. 35]
\textbf{Алгеброй над полем} $K$ называется множество $A$ с операциями сложения, 
умножения и умножения на элементы поля $K$, обладающими следующими свойствами:
\begin{enumerate}
    \item[1)] относительно сложения и умножения на элементы поля $A$ есть векторное пространство;
    \item[2)] относительно сложения и умножения $A$ есть кольцо;
    \item[3)] $(\lambda a)b = a(\lambda b) = \lambda(ab)$ для любых $\lambda \in K$, $a, b \in A$.
\end{enumerate}

\begin{example}
Всякое поле $L$, содержащее $K$ в качестве подполя, можно 
рассматривать как алгебру над $K$. В частности, поле $\mathbb{C}$ есть алгебра над $\mathbb{R}$.
\end{example}

\begin{example}
Пространство $E^3$ есть алгебра относительно операции векторного умножения.
\end{example}

\end{definition}

\subsection*{2. Линейная оболочка и способы задания подпространств}

\begin{definition}
Пусть в линейном пространстве $V$ задана система векторов $a_1, \dots, a_k$ (необязательно линейно независимых). Тогда множество $\mathcal{L}(a_1, \dots, a_k) = \{ \lambda_1 a_1 + \dots + \lambda_k a_k \mid \lambda_i \in \mathbb{F} \}$ (то есть все возможные линейные комбинации векторов $a_1, \dots, a_k$) называется \textbf{линейной оболочкой} системы векторов $a_1, \dots, a_k$.
\end{definition}

\subsubsection*{Задание подпространства}

Существует два основных способа задания подпространства $L$:
\begin{enumerate}
    \item \textbf{Линейная оболочка:} $L = \langle e_1, \dots, e_k \rangle$
    \item \textbf{Через ОСЛАУ:} $L = \{ x \in \mathbb{R}^n \mid Ax = 0 \}$
\end{enumerate}

\subsubsection*{Связь между способами задания}

Путь $L = \langle e_1, \dots, e_k \rangle$. Переход к ОСЛАУ строится по следующей цепочке эквивалентностей:

\begin{align*}
x \in L &\iff Ax = 0 \\
&\iff \forall e_i : Ae_i = 0 \\
&\iff (e_i)^T \cdot A^T = 0 \\
&\iff 
\begin{pmatrix}
e_1^T \\
e_2^T \\
\vdots \\
e_k^T
\end{pmatrix} \cdot A^T = 0 \\
&\iff \text{столбцы } A^T \text{ (строки } A\text{) являются решением для ОСЛАУ: } \\
& \quad \begin{pmatrix}
e_1^T \\
e_2^T \\
\vdots \\
e_k^T
\end{pmatrix} y = 0
\end{align*}

\noindent\textbf{Пошаговый практический алгоритм}

Чтобы перевести линейную оболочку $\langle e_1, \dots, e_k \rangle$ в ОСЛАУ $Ax = 0$:

\begin{enumerate}
    \item \textbf{Запишите матрицу $E$:} положите данные векторы $e_i$ в строки матрицы.
    \item \textbf{Составьте вспомогательную ОСЛАУ:} $Ey = 0$.
    \item \textbf{Найдите её фундаментальную систему решений (ФСР):} решите эту систему методом Гаусса.
    \item \textbf{Сформируйте матрицу $A$:} полученные векторы ФСР запишите как строки искомой матрицы $A$.
\end{enumerate}

Система $Ax = 0$ и будет искомой ОСЛАУ для вашего подпространства.


\subsection*{3. Сумма и пересечение подпространств}

\begin{definition} Суммой подпространств $L_1$ и $L_2$ пространства $L$ называется множество
$$L_1 + L_2 = \{x \in L \mid \exists x_1 \in L_1, x_2 \in L_2 : x = x_1 + x_2\}.$$
\end{definition}

\begin{definition} Пересечением подпространств $L_1$ и $L_2$ пространства $L$ называется множество
$$L_1 \cap L_2 = \{x \in L \mid x \in L_1 \wedge x \in L_2\}.$$
\end{definition}

\begin{definition} Сумма подпространств $L_1$ и $L_2$ пространства $L$ называется \textit{прямой суммой} и обозначается $L_1 \oplus L_2$, если $\dim(L_1 \cap L_2) = 0$.
\end{definition}

\begin{theorem}
$\dim L_1 + \dim L_2 = \dim(L_1 + L_2) + \dim(L_1 \cap L_2)$.
\end{theorem}
